\section{Conclusion}
\label{sec:conclusion}

In this paper, we presented a pipeline for the discovery of novel biological mechanisms through the lens of mechanistic interpretability. By combining the disentanglement capabilities of Sparse Autoencoders with the semantic reasoning of Large Language Models, we identified and characterized several "orphan" features in the \esm protein language model. Our findings suggest that these models learn highly specific biophysical abstractions, such as charged turn motifs and helix stabilizers, which are not currently captured in standard biological annotations.

The ability to bridge the gap between high-dimensional neural activations and human-readable biological hypotheses represents a significant step toward using AI as an active partner in scientific discovery. While our study is a proof-of-concept, it highlights the potential for foundation models to reveal deeper principles of biological organization.

Future work will focus on scaling this pipeline to the entire \swissprot database and implementing automated filters to exclude features that correlate with known experimental structures (e.g., from the PDB). Most importantly, we aim to validate the generated hypotheses through laboratory experiments, testing whether the identified "orphan" features indeed correspond to critical determinants of protein function and stability.
